%! Setting Up --- Learning from Data — Unit 8, Lesson 8.0 — Lesson Plan
\documentclass[10pt]{article}
\usepackage{linalg-article}
\usepackage{linalg-boxes}
\usepackage{graphicx}   % -article does NOT load graphicx; needed for thumbnails/screenshots
\graphicspath{{images/}}
\newcommand{\TallMath}[1]{$\displaystyle #1\rule[-1.4em]{0pt}{3.2em}$}
\newcommand{\vv}{\mathbf{v}}
\newcommand{\ww}{\mathbf{w}}
\newcommand{\bb}{\mathbf{b}}
\newcommand{\zero}{\mathbf{0}}
\newcommand{\relu}{\mathrm{ReLU}}
\newcommand{\UnitNumberName}{Unit 8: Learning from Data \quad}
\newcommand{\LessonNumberName}{Lesson 8.0: Setting Up --- Learning from Data}

\begin{document}
\begin{center}
    {\Large\bfseries \CourseName: \SchoolYear \\
    {\small\bfseries \UnitNumberName \LessonNumberName}}
\end{center}

\begin{tcolorbox}[colback=blush, colframe=burgundy, boxrule=0.9pt, arc=2mm,
    left=2mm, right=2mm, top=1.5mm, bottom=1.5mm]
    \textbf{Primary Objective:} Students can describe a \emph{learning function} (a neural network) as a
    chain of \emph{layers}, each doing a \emph{linear step} --- multiply by a weight matrix and add a bias
    (Unit~1) --- then the one \emph{nonlinear step}, the \emph{ReLU} $=\max(0,x)$. They can evaluate ReLU
    and compute one layer $\relu(A\vv+\bb)$, and explain how stacking and shifting ReLUs builds
    \emph{piecewise-linear} functions that bend to fit data (unlike a single straight line). They can
    describe \emph{learning} as choosing the weights to shrink a \emph{loss} (the Unit~4 sum of squared
    errors) by \emph{gradient descent}, and articulate the unit's big idea: modern AI is, at heart, the
    linear algebra of this course. They can preview where Unit~8 goes --- piecewise-linear functions (8.1),
    convolutional nets (8.2), gradient descent (8.3), and mean/variance/covariance (8.4).
    \emph{(Unit~8 is optional enrichment --- explored, not formally assessed.)}
\end{tcolorbox}

\renewcommand{\arraystretch}{1.4}
\begin{skillbox}[Priority Ideas \& Skills]{goldbox}
\begin{tabularx}{\linewidth}{@{}X|X@{}}
\textbf{Priority Ideas \& Skills} & \textbf{Key Understandings (the ``why'')} \\[2pt]
\hline
\vspace{-6pt}
\begin{itemize}[leftmargin=1.1em, nosep, after=\vspace{2pt}]
  \item Evaluate the \textbf{ReLU} $=\max(0,x)$; know it is one \textbf{fold}.
  \item Compute one \textbf{layer} $\relu(A\vv+\bb)$ (linear step, then bend).
  \item Explain how stacked ReLUs make \textbf{piecewise-linear} functions that bend to fit data.
  \item Read a \textbf{loss} (sum of squared errors); describe learning as \textbf{gradient descent}.
\end{itemize}
&
\vspace{-6pt}
\begin{itemize}[leftmargin=1.1em, nosep, after=\vspace{2pt}]
  \item A network is \emph{matrix multiplication} (Unit~1) $+$ one simple bend --- no new machinery.
  \item The ReLU's nonlinearity is what keeps stacked layers from collapsing into a single matrix.
  \item ``Learning'' $=$ choosing weights to make the least-squares loss (Unit~4) as small as possible.
  \item The linear algebra of this course is the engine inside modern AI.
\end{itemize}
\end{tabularx}
\end{skillbox}

\begin{skillbox}[{Vocabulary, Concepts, \& Theorems}]{greenbox}
\renewcommand{\arraystretch}{1.3}
\begin{tabularx}{\linewidth}{@{}l X@{}}
\textbf{Learning function $F(\vv)$} & A neural network: turns an input vector into a prediction by alternating linear steps with ReLUs. \\
\textbf{ReLU $=\max(0,x)$} & The nonlinear step: keeps a positive number, replaces a negative with $0$. One fold. \\
\textbf{Weight matrix $A$, bias $\bb$} & The numbers a network chooses; $A$ scales/mixes the inputs, $\bb$ shifts them. \\
\textbf{One layer $\relu(A\vv+\bb)$} & Linear step ($A\vv+\bb$, Unit~1) then ReLU applied to each entry. \\
\textbf{Piecewise-linear function} & Straight pieces joined at folds; sums of ReLUs build them, and they can bend to fit any data. \\
\textbf{Loss} & Sum of squared errors between predictions and targets (Unit~4 least squares). \\
\textbf{Gradient descent} (preview 8.3) & Nudge the weights ``downhill'' to make the loss smaller. \\
\end{tabularx}
\end{skillbox}

\begin{fixedskillbox}[Activate Prior Knowledge \& Spiral Review]{blush}
    The Warm-Up rehearses the three pieces a neural network is built from:
    \textbf{(1)} a \emph{matrix times a vector} $A\vv$ (Unit~1) --- the linear step of every layer;
    \textbf{(2)} \emph{squared error} (Unit~4 least squares) --- a guess minus a target, squared, is how a
    network scores ``how wrong am I''; and \textbf{(3)} the plain-arithmetic \emph{$\max(0,x)$ function} ---
    which is exactly the \textbf{ReLU}, the one new idea today. Debrief item~3 aloud: $\max(0,x)$ keeps
    positives and zeroes negatives, so it is the ``bend'' the whole unit rides on.
\end{fixedskillbox}

\begin{skillbox}[Hook]{blush}
    Open with the punchline: the matrices we have used all year to move, project, stretch, and factor the
    plane are the \emph{engine inside modern AI}. A \textbf{neural network} is just a \emph{function} that
    turns an input into a prediction --- and it is built almost entirely from Unit~1 (matrix times vector)
    and Unit~4 (least squares), plus \emph{one} new simple function. Ask: \textbf{``The Unit~4 best-fit
    line has two weights (slope, intercept) --- what can it \emph{not} do?''} Answer: bend. Name the payoff
    question of Unit~8: \textbf{``How do a matrix and one simple bend combine into a machine that learns
    from data?''} That bend is the ReLU; stacking it builds functions that fit anything; and ``learning''
    is just shrinking a loss. Note for students: Unit~8 is a \emph{tour} --- optional enrichment we explore
    for the big picture, not a formally tested unit.
\end{skillbox}

\begin{skillbox}[Lesson]{blush}
\begin{multicols}{2}
\textbf{1. Learning $=$ fitting a function.} Given input--output examples, find a function that reproduces
them. The simplest is the Unit~4 best-fit line $y=wx+b$ (two weights). A line only goes straight; real
data \emph{bends}, so we need functions that bend while still built from straight pieces.

\textbf{2. The one new ingredient: ReLU.} $\relu(x)=\max(0,x)$ keeps positives, zeroes negatives ---
one \textbf{fold} at the origin. Add/shift a few ReLUs and the folds combine into a
\textbf{piecewise-linear} function that bends to fit data. E.g.\ $\relu(x)-\relu(x-2)$ ramps up then
levels off (two folds). This is Lesson~8.1.

\columnbreak

\textbf{3. A layer is a matrix step.} One layer maps $\vv\mapsto\relu(A\vv+\bb)$: linear step (matrix
$\times$ vector $+$ bias, Unit~1) then ReLU on each entry. On $A=\begin{bsmallmatrix} 1 & -1\\ 1 & 1
\end{bsmallmatrix}$, $\vv=\begin{bsmallmatrix} 1\\2 \end{bsmallmatrix}$: $A\vv=\begin{bsmallmatrix} -1\\3
\end{bsmallmatrix}\to\relu\to\begin{bsmallmatrix} 0\\3 \end{bsmallmatrix}$. Matrix multiplication does the
heavy lifting; $\bb$ shifts the folds.

\textbf{4. Learning + why it matters.} Score predictions with a \textbf{loss} $=$ sum of squared errors
(Unit~4); \textbf{learn} by nudging weights downhill (\emph{gradient descent}, 8.3). A network is Unit~1
$+$ one bend $+$ Unit~4 --- no new machinery. \textbf{Road ahead:} 8.1 piecewise-linear; 8.2 conv nets;
8.3 gradient descent; 8.4 mean/variance/covariance.
\end{multicols}
\end{skillbox}

\begin{skillbox}[Explicit Instruction: Running One Layer $\relu(A\vv+\bb)$]{blush}
\begin{tabularx}{\linewidth}{@{}p{0.46\linewidth} X@{}}
\textbf{Steps (one layer).}
\begin{enumerate}[leftmargin=1.3em, nosep, after=\vspace{2pt}]
  \item Linear step: compute $A\vv$ (row $\cdot$ vector, Unit~1).
  \item Add the bias: $A\vv+\bb$ (shifts each entry).
  \item Apply ReLU to \emph{each entry}: keep positives, zero negatives.
  \item The result is the layer's output --- it feeds the next layer.
\end{enumerate}
&
\textbf{Worked example} ($A=\begin{bsmallmatrix} 1 & -1\\ 1 & 1 \end{bsmallmatrix}$, $\bb=\zero$,
$\vv=\begin{bsmallmatrix} 1\\2 \end{bsmallmatrix}$). Linear step:
$A\vv=\begin{bsmallmatrix} -1\\3 \end{bsmallmatrix}$. ReLU each entry: $-1\to0$, $3\to3$, giving
$\begin{bsmallmatrix} 0\\3 \end{bsmallmatrix}$. \textbf{Common slip:} applying ReLU to the vector as a
whole, or ``fixing'' the positive $3$. (Bias check: with $\bb=\begin{bsmallmatrix} -1\\0 \end{bsmallmatrix}$
the first entry becomes $-2\to0$ --- the bias moved the fold.)
\end{tabularx}
\end{skillbox}

\begin{skillbox}[Active Monitoring]{redbox}
Circulate during the notes practice and Tier work. Watch for and correct:
\begin{itemize}[leftmargin=1.2em, nosep]
  \item arithmetic slips in $A\vv$ (row $\cdot$ vector) --- every layer starts with the Unit~1 product;
  \item applying ReLU to the \emph{whole} vector instead of \emph{each entry};
  \item thinking ReLU changes positive numbers (it only touches negatives);
  \item forgetting to square (or dropping the sign before squaring) when computing the loss.
\end{itemize}
Cold-call prompts: ``What does ReLU do to $-1$? to $5$?'' ``Which entry did the ReLU change?'' ``Is a
lower loss better or worse?'' ``Without a ReLU, what do two matrix layers collapse into?''
\end{skillbox}

\begin{skillbox}[Group Work \& Differentiation]{redbox}
\textbf{Building a Learning Function} activity (tiered; mirrors the notes).
\begin{multicols}{3}
\textbf{Tier R --- Remediate.} Evaluate ReLU; run one linear step $A\vv$ then ReLU, and say which entry
the ReLU changed.

\columnbreak
\textbf{Tier A --- Approaching.} A full layer with a bias $\relu(A\vv+\bb)$; a bending function
$\relu(x)+\relu(x-2)$ (find its folds \& slopes); and comparing two networks by their loss (which fits
better, which way learning moves).

\columnbreak
\textbf{Tier E --- Extension.} \emph{Why the ReLU is not optional:} two linear layers with no ReLU
collapse into one matrix $A_2A_1$ (a single straight-line map) --- the ReLU's fold is what makes a deep
network more than one matrix.
\end{multicols}
\end{skillbox}

\begin{skillbox}[Individual Work \& Assessment]{redbox}
\textbf{Exit Ticket} (independent, no notes): compute one layer $\relu(A\vv)$; compute a loss (sum of
squared errors); and a \textbf{conceptual/justification check} --- in one sentence, why a network needs
the nonlinear ReLU (else stacked layers collapse into a single matrix that cannot bend). Sort into ``got
it / arithmetic slip / concept slip'' to decide whether 8.1 opens with a quick ``one ReLU $=$ one fold''
recap. \emph{Enrichment unit --- use as formative feedback, not a graded assessment.}
\end{skillbox}

\begin{skillbox}[Reinforcement \& Extension]{goldbox}
\textbf{Homework:} evaluate ReLU (incl.\ a decimal and $\relu(0)$); run a layer with a bias on two inputs;
evaluate a piecewise-linear $\relu(x)-\relu(x-3)$ and locate its folds; compare two networks by loss; and
short justifications (matrix multiply as the engine; why the ReLU is essential). \textbf{Extension:}
compute $A_2A_1$ to see two ReLU-less layers collapse into one matrix. \textbf{Preview:} Lesson~8.1,
\emph{Piecewise Linear Learning Functions} --- folding many lines with ReLUs into a continuous
piecewise-linear function that can fit any data.
\end{skillbox}
\end{document}
